\scanpage{16}
Summary of some results of Kostant (simple subgroups of rank $1$).

$G$ is a semi-simple group over a field $k$ of characteristic zero, $\handwritten{\mathfrak{g}}$ its Lie algebra,
\margin{\handwritten{(The principal 3-dimensional subgroup...., Am. Jour. Oct.1959, 3.4, 3.5, 3.6, \unsure{1.2|4.2}, 9.3)}}

\uline{Theorem} 1. \inserted{\handwritten{(i)}} Le\deleted{\unsure{t}}t $x,e\handwritten{\in\mathfrak{g}}$,\inserted{with $e\handwritten{\ne}0$, $e$ nilpotent.} For every integer $k$, let $\handwritten{\mathfrak{g}}(k)$ be the eigenspace of\hspace{1.5em} corresponding to the eigenvalue $k$. Thus if $x=\operatorname{Lie}(i)(1)$, where $i:\underline{G}_m\handwritten{\longrightarrow}G$, the $\handwritten{\mathfrak{g}}(k)$ are the \deleted{decomposition} components of $\handwritten{\mathfrak{g}}$ according to the characters of $\underline{G}_m$, acting on $\handwritten{\mathfrak{g}}$ via $i$ and the adjoint representation. The following conditions are equiva\deleted{\unsure{e}}lent:

\[
\text{a)}\quad e\handwritten{\in\mathfrak{g}}(2)\quad\text{i.e.}\quad \handwritten{[}x,e\handwritten{]}=2e,\quad\text{and }\inserted{\text{for every integer }k\handwritten{\geq}1,}\quad\operatorname{ad}(e)^k:\handwritten{\mathfrak{g}}(-k)\handwritten{\longrightarrow\mathfrak{g}}(k)\quad\text{is an isomorphism.}
\]

\[
\text{b)}\quad e\handwritten{\in\mathfrak{g}}(2),\quad\handwritten{\deleted{\unsure{and}}}\quad \operatorname{ad}(e):\handwritten{\mathfrak{g}}(0)\handwritten{\longrightarrow\mathfrak{g}}(-2)\quad\text{is surjective},
\]
\inserted{\handwritten{and $x$ is of the form $f(x_o)$, for $f:\operatorname{Sl}(2,k)\to G$ (\uline{Dynkin element}).}}

\handwritten{\deleted{b') As in b), with $\operatorname{ad}(e):\mathfrak{g}(0)\to\mathfrak{g}(\unsure{2|-2})$ \uline{\unsure{bijective}}}}

\[
\text{c)}\quad e\handwritten{\in\mathfrak{g}}(2),\quad\text{and }x\handwritten{\in}\operatorname{ad}(e)(\handwritten{\mathfrak{g}}).
\]

\[
\text{c')}\quad e\handwritten{\in\mathfrak{g}}(2),\quad\text{and }x\handwritten{\in}\operatorname{ad}(e)(\handwritten{\mathfrak{g}}(-2)).
\]

d) There exists a homomorphism $f:\operatorname{Sl}(2,k)\handwritten{\longrightarrow}G$ such that $f^{\handwritten{'}}(e_o)=e,f^{\handwritten{'}}(x_o)=x\deleted{y}$ where $(x_o,e_o,f_o)$ is the basis of $\operatorname{Lie}\operatorname{Sl}(2,k)$ corresponding to the usual pinning of this group.

Th$\overset{-}{\underset{\text{\unsure{e}}}{\text{e}}}$se conditions imply that the\deleted{\unsure{o|b}}re exists a homomorphism $i:\underline{G}_m\handwritten{\longrightarrow}G$, such that $\operatorname{Lie}(i)(1)=x$, a fortiori that $x$ is semi-simple. Moreover, the $f$ of d) is uniquely determined.

(ii) For every nilpotent element $e$ of $\handwritten{\mathfrak{g}}$, there exists a homomorphism $f:\operatorname{Sl}(2,k)\handwritten{\longrightarrow}G$ such that $f(e_o)=e$, or what amounts to the same thing, such that $f(u_o)=u$ (where $u=\exp(e)$, $u_o=\exp(e_o)$);\inserted{\handwritten{$f$ is uniquely determined \unsure{by} $x=f(x_o)$ i.e., by $i=f\circ i_o$.}} \deleted{\unsure{e}}\handwritten{B}y virtue of (i), it amounts to the same thing,\deleted{\unsure{(}}when $e=0$ i\handwritten{/}.e. $u\handwritten{\ne}1$, to say that there exists an $x\handwritten{\in\mathfrak{g}}$ such that $(x,e)$ satisfies the equivalent conditions of (i), or again a homomorphism $i:\underline{G}_m\handwritten{\longrightarrow}G$ such that $(\operatorname{Lie}(i)(1),e)$ satisfies them. The homomorphisms $f$ in question are conjugate \deleted{under the elements of}\inserted{by $\handwritten{Z=}$}$\operatorname{Cent}(e)=\operatorname{Cent}(u)$ (i.e. for $e\handwritten{\ne}0$, the $x$ in question, or again the $i$ in question, are conjugate \deleted{under} by $\operatorname{Cent}(u)$), and when $e\handwritten{\ne}0$, the \deleted{\unsure{morphisms}\quad $g$\quad $\operatorname{int}(g)$} $f$ in question (or again the
% archive sheet number: 15
